% ============================================
% DOC-21 (ZH)
% 中文翻译 — SPEC-935-R356
% ============================================
\section*{\textbf{DOC-21} --- 埃托尔·奥利维拉的信（未寄出）}
\markboth{DOC-21 --- 埃托尔·奥利维拉的信（未寄出）}{}

\noindent \textit{从鉴定档案中提取的文件。}

\vspace{0.5em}\noindent\rule{\textwidth}{0.5pt}\vspace{0.5em}


\noindent {\footnotesize\texttt{公共档案馆 $\cdot$ 书证 $\cdot$ 案卷1,260}}\\[-0.2em]

\noindent \textbf{收信人:} 玛丽亚·奥利维拉（妻子）。\\
\noindent \textbf{寄信人:} 埃托尔·奥利维拉医生。\\
\noindent \textbf{地点:} 巴尔巴塞纳。\\
\noindent\textbf{日期:} 未登记（1979年6月16日至6月20日之间）。\\
\noindent\textbf{原件状态:} 在奥利维拉医生私人笔记本的封底内侧被发现，折成四折，没有信封。未寄出。

\noindent 玛丽亚，

\noindent 如果你正在读这封信，我已经离开了家。或者离开了我自己。我不知道
\noindent 哪个先发生。

\noindent 1,260的日记在我这里。我没有请求许可就把它从法医研究所带走了。
\noindent 我在登记簿上签了字，装作有授权，没有人质疑。这里就是这样运作的：医生一签
\noindent 字，没有人问。

\noindent 我以为两天就能解决。拿走笔记本，从容地读，写一份意见书，还回去。例行公事。
\noindent 那个人已经死了六十二年。还能发生什么？

\noindent 玛丽亚，我没有读那本日记。是日记读了我。

\noindent 第三天，我醒来时眼睛里有一块黄斑。不疼。
\noindent 不痒。但它正在扩大。眼科医生说这是结膜下出血。两周内会消失。但我知道那不是
\noindent 血。那是患者在40页描述的眼睛的颜色。我今天早上重读了40页。描述和我眼睛里的
\noindent 斑块一模一样。

\noindent 第五天，我闻到了那股气味。

\noindent 你认识收容院的气味。1975年你跟我来过一次。
\noindent 你说再也不来了。记得吗？来苏水、粪便、回锅饭、密闭身体的气味。日记有那种气味。但不是来自收容院。来自纸张。仿佛墨水是用某种腐烂后又被遗忘的东西制成的。

\noindent 第七天，我开始书写。

\noindent 不是我在写。是手。它自己动，和患者描述的一样。我看着字迹出现——不是我的。是他的。同样的左倾。同样的开口“a”。同样的不规律压力，仿佛笔太重了。

\noindent 今天我醒来，不知道自己在哪里。我在一段不属于我的记忆里醒来。一道铁丝网围栏。太阳灼烧着后颈。太瘦的人靠着铁丝，眼睛一动不动。我从未见过那景象。但身体记得。双手颤抖，仿佛握过滚烫的铁丝。

\noindent 玛丽亚，我害怕睡觉。每次闭上眼睛，我就回到乱葬坑。就是他描述的那个坑。同样的手臂从土里伸出。同样的甜味粘在鼻子里出不来。

\noindent 我不会把日记还回去。

\noindent 我知道这听起来像疯了。但如果我还回去，就会有人读。而如果有人读，同样的事情就会发生。我不能对任何人做这件事。

\noindent 或者——我不知道自己是在保护别人，还是日记不想被还回去。

\noindent 它喜欢我。

\noindent 它喜欢和我在一起。

\noindent 它不想回抽屉。

\noindent 对不起。

\noindent 埃托尔

\noindent 档案员注：这封信于2026年被发现，在奥利维拉医生的私人笔记本内，与患者1,260的日记一起被归还给米纳斯吉拉斯州公共档案馆，距信写成已47年。纸上的折痕表明，他在这段时间里一直随身携带着它。





\vspace{0.5em}
\noindent\textit{\textbf{↪ 路线:} \rota{DOC-01} \rotasep \rota{LUC-02} \rotasep \rota{DOC-22}}
